\section*{الفصل \ref{ch:g12:ka-and-pka} --- الأحماض والقواعد: Ka وpKa}

\begin{solution}{exo:g12:ka-and-pka:1}
القواعد المرافقة: \ce{HCOO-}، \ce{NH3}، \ce{OH-}، \ce{CO3^{2-}}. والأحماض
المرافقة: \ce{NH4+}، \ce{H2O}، \ce{H3O+}، \ce{CO2,H2O}. والأمفوليتات: \ce{H2O}
و\ce{HCO3-}.
\end{solution}

\begin{solution}{exo:g12:ka-and-pka:2}
\omterm{def:g9:acids-bases-ph:ph}{pH} 1.0؛ 3.0؛ $-\log(\num{5.0e-3}) = 2.30$.
\end{solution}

\begin{solution}{exo:g12:ka-and-pka:3}
$10^{-4.5} = \qty{3.2e-5}{mol/L}$. وعند \omterm{def:g9:acids-bases-ph:ph}{pH} 11.2:
$[\ce{H3O+}] = \qty{6.3e-12}{mol/L}$
و$[\ce{OH-}] = \num{1.0e-14}/\num{6.3e-12} = \qty{1.6e-3}{mol/L}$.
\end{solution}

\begin{solution}{exo:g12:ka-and-pka:4}
\ce{HCOOH + H2O <=> HCOO- + H3O+}،
$K_a = \dfrac{[\ce{HCOO-}][\ce{H3O+}]}{[\ce{HCOOH}]\,c^\circ}$؛
\ce{NH4+ + H2O <=> NH3 + H3O+}،
$K_a = \dfrac{[\ce{NH3}][\ce{H3O+}]}{[\ce{NH4+}]\,c^\circ}$.
\end{solution}

\begin{solution}{exo:g12:ka-and-pka:5}
حمض الميثانويك ($pK_a$ أصغر). والقاعدة الأقوى هي \ce{CH3COO-}،
قاعدة الحمض الأضعف.
\end{solution}

\begin{solution}{exo:g12:ka-and-pka:6}
$[\ce{A-}] = [\ce{H3O+}] = \qty{1.0e-3}{mol/L}$؛
$[\ce{HA}] = 0.050 - 0.001 = \qty{0.049}{mol/L}$؛
$K_a = \num{1.0e-6}/0.049 = \num{2.0e-5}$؛ $pK_a = 4.69$.
\end{solution}

\begin{solution}{exo:g12:ka-and-pka:7}
$K_a = \num{6.31e-5}$؛ $x^2 + \num{6.31e-5}\,x - \num{1.26e-6} = 0$ تعطي
$x = \qty{1.09e-3}{mol/L}$: \omterm{def:g9:acids-bases-ph:ph}{pH} 2.96.
\end{solution}

\begin{solution}{exo:g12:ka-and-pka:8}
$\mathrm{pH} = 14.0 + \log(\num{2.0e-3}) = 14.0 - 2.70 = 11.3$.
\end{solution}

\begin{solution}{exo:g12:ka-and-pka:9}
\ce{HF}، \ce{HCOOH}، حمض الأسكوربيك، \ce{C6H5COOH}، \ce{CH3COOH}،
\ce{CO2,H2O}، \ce{NH4+}، \ce{HCO3-}. ويقع \omterm{def:g12:ka-and-pka:strong-acid}{الحمض القوي} في الأسفل بعيدًا،
وراء طرف السلّم: فهو يتفاعل مع الماء \omterm{def:g10:reaction-progress-table:limiting}{تفاعلًا تامًا}.
\end{solution}

\begin{solution}{exo:g12:ka-and-pka:10}
$[\ce{H3O+}] = 10^{-5.6} = \qty{2.5e-6}{mol/L}$، أي أكثر بخمس وعشرين مرة مما في
الماء النقي (\qty{1.0e-7}{mol/L}).
\end{solution}

\begin{solution}{exo:g12:ka-and-pka:11}
$K_a K_b = \dfrac{[\ce{NH3}][\ce{H3O+}]}{[\ce{NH4+}]c^\circ} \times
\dfrac{[\ce{NH4+}][\ce{OH-}]}{[\ce{NH3}]c^\circ} = K_e$، إذن
$pK_a = pK_e - pK_b = 14.0 - 4.74 = 9.26$.
\end{solution}

\begin{solution}{exo:g12:ka-and-pka:12}
عند \qty{0.0010}{mol/L}: $x^2 + \num{1.74e-5}\,x - \num{1.74e-8} = 0$،
$x = \qty{1.24e-4}{mol/L}$، \omterm{def:g9:acids-bases-ph:ph}{pH} 3.91: ارتفاع قدره 0.52، لا 1. \omterm{def:g12:ka-and-pka:strong-acid}{فالحمض الضعيف} إذا
خُفّف تفاعل مع الماء بنسبة أكبر (هنا
\qty{12}{\percent} مقابل \qty{4}{\percent})، وهذا يعوّض جزئيًا
\omterm{def:g10:concentration-and-dilution:dilution}{التخفيف}.
\end{solution}

\begin{solution}{exo:g12:ka-and-pka:13}
بوصفه حمضًا: \ce{HCO3- + H2O <=> CO3^{2-} + H3O+}، $K = 10^{-10.33}$. وبوصفه
قاعدة: \ce{HCO3- + H2O <=> H2CO3 + OH-} (حيث يمثل \ce{H2CO3} الثنائي \ce{CO2,H2O})، $K = K_e/K_a = 10^{-14.0+6.37}
= 10^{-7.63}$. ولتفاعل القاعدة الثابت الأكبر: \omterm{def:g9:acids-bases-ph:acidic}{فالمحلول
قاعدي} قليلًا.
\end{solution}

\begin{solution}{exo:g12:ka-and-pka:14}
لا يستطيع حمض، مهما خُفّف، أن يجعل \omterm{def:g2:mixing-and-dissolving:dissolve}{المحلول} قاعديًا. فحين تصير \omterm{def:g9:ions:ion}{أيونات}
الأكسونيوم الآتية من الحمض أقل من \qty{1e-7}{mol/L} التي يوفرها
الماء، تغلب \omterm{def:g9:ions:ion}{أيونات} الماء: فتقترب \omterm{def:g9:acids-bases-ph:ph}{pH} من 7 من الأسفل
ولا تتجاوزها أبدًا.
\end{solution}

\begin{solution}{exo:g12:ka-and-pka:15}
$n = 0.500 / 176.0 = \qty{2.84e-3}{mol}$، $c = \qty{0.0142}{mol/L}$؛
$K_a = \num{9.12e-5}$؛ $x^2 + \num{9.12e-5}\,x - \num{1.30e-6} = 0$ تعطي
$x = \qty{1.09e-3}{mol/L}$: \omterm{def:g9:acids-bases-ph:ph}{pH} 2.96.
\end{solution}

\begin{solution}{pb:g12:ka-and-pka:1}
\textbf{1.} \ce{H-C(=O)-O-H}: والهيدروجين المعطى هو هيدروجين \ce{O-H}.

\textbf{2.} $\ce{HCOOH}/\ce{HCOO-}$؛ \ce{HCOOH + H2O <=> HCOO- + H3O+}.

\textbf{3.} $K_a = \dfrac{[\ce{HCOO-}][\ce{H3O+}]}{[\ce{HCOOH}]c^\circ} =
10^{-3.74} = \num{1.82e-4}$.

\textbf{4.} تحت حمض الإيثانويك ($3.74 < 4.76$): فهو أقوى.

\textbf{5.} $M = \qty{46.0}{g/mol}$: $0.010 \times 46.0 = \qty{0.46}{g}$.

\textbf{6.} \ce{HCOOH}: $0.010 - x$؛ \ce{HCOO-}: $x$؛ \ce{H3O+}: $x$.

\textbf{7.} $\dfrac{x^2}{0.010 - x} = \num{1.82e-4}$.

\textbf{8.} $x^2 + \num{1.82e-4}\,x - \num{1.82e-6} = 0$:
$x = \qty{1.26e-3}{mol/L}$.

\textbf{9.} $\mathrm{pH} = -\log(\num{1.26e-3}) = 2.90$.

\textbf{10.} $\num{1.26e-3}/0.010 = \qty{12.6}{\percent}$.

\textbf{11.} \qty{1.26e-3}{mol/L} أكبر من \qty{1e-7}{mol/L} بأكثر من عشرة آلاف
مرة: \omterm{def:g9:ions:ion}{فأيونات} الماء نفسه مهملة.

\textbf{12.} \omterm{def:g9:acids-bases-ph:ph}{pH} 2.0.

\textbf{13.} $0.010 / \num{1.26e-3} \approx 8$ مرات أكثر.

\textbf{14.} حمض الهيدروكلوريك \omterm{def:g12:ka-and-pka:strong-acid}{حمض قوي}، يتحول كله إلى
\omterm{def:g9:ions:ion}{أيونات}؛ وحمض الميثانويك ضعيف: فتفاعله مع الماء يتوقف عند
توازن يبقى فيه معظمه على هيئة \ce{HCOOH}.

\textbf{15.} \ce{HCOOH + HCO3- -> HCOO- + CO2 + H2O}.

\textbf{16.} بضرب كلٍّ من بسط الكسر الآتي ومقامه
\[
  K = \frac{[\ce{HCOO-}]\,[\ce{CO2}]}{[\ce{HCOOH}]\,[\ce{HCO3-}]}
\]
في $[\ce{H3O+}]$ نحصل على $K = K_a(\ce{HCOOH}) / K_a(\ce{CO2,H2O}) =
\num{1.82e-4}/\num{4.3e-7} \approx 420$.

\textbf{17.} نعم، $K \gg 1$؛ ويتصاعد ثنائي أكسيد الكربون فقاعات.

\textbf{18.} الهيدروكسيد \omterm{def:g12:ka-and-pka:strong-acid}{قاعدة قوية} تهاجم الجلد هي نفسها؛
أما الهيدروجينوكربونات \omterm{def:g12:ka-and-pka:strong-acid}{فقاعدة ضعيفة} تعادل الحمض وتبقى لطيفة
حتى إذا زادت عن الحاجة.

\textbf{19.} \omterm{def:g9:acids-bases-ph:ph}{pH} 2.9.
\end{solution}
